% Integration-ready research fragment. No custom macros are required.
% The ambient article should provide theorem, proposition, corollary, proof.
\section{Integer dilations and exact local scale corrections}
\label{dil:section}

Two unequal integer frequencies still leave a single free Fourier index.
This elementary fact resolves the unequal-dilation question in the incoming
correlation reports. The spectral reduction is classical in method:
Espinosa--Moll, Part~1, equation~(1.13), already records an unshifted,
equal-order dilation formula due to Mikol\'as. The additional issue here is
the finite part in an explicitly fixed local coordinate, together with all
spectral and argument derivatives. No historical priority claim is made for
the underlying Fourier factorization.

\subsection{The reduced frequency and shifted spectral kernel}
Let $p,q\geq1$ be integers, let
\[
 d=\gcd(p,q),\qquad p=dP,\quad q=dQ,\quad \gcd(P,Q)=1,
 \qquad c=\{Pb-Qa\}.
\]
The shifts $a,b$ are real and henceforth $0<c<1$. This condition is exactly
that the singular sets of the two periodic factors are disjoint. Indeed,
simultaneous singularities would give $px+a=j$, $qx+b=k$ and hence
$Pb-Qa=Pk-Qj\in\mathbb Z$; conversely, Bezout's identity gives such a
solution whenever this expression is integral.

Write
\[
 C(s,t;c)=\int_0^1\zeta(s,x)\zeta(t,\{x+c\})\,dx,
 \qquad \Re s,\Re t<1.
\]

\begin{theorem}[Complete bilinear integer-dilation reduction]
\label{dil:master}
For $\Re s,\Re t<1$,
\begin{equation}\label{dil:factor}
 \int_0^1\zeta(s,\{px+a\})\zeta(t,\{qx+b\})\,dx
 =Q^{s-1}P^{t-1}C(s,t;c).
\end{equation}
Both sides are jointly holomorphic in this domain and their meromorphic
continuations agree. Thus arbitrary spectral derivatives follow by
differentiating the displayed formula.
\end{theorem}

\begin{proof}
The disjoint endpoint singularities imply local domination by
$t^{-\sigma}(1+|\log t|^M)$ for some $\sigma<1$ on every compact parameter
set. This proves convergence and holomorphy, including all spectral jets.
Initially take $\Re s,\Re t<0$, where the Fourier series converge
absolutely. A surviving frequency satisfies $pn+qm=0$, whose complete
integer solution is $n=Qr$, $m=-Pr$. For nonzero $r$, the Hurwitz Fourier
coefficient satisfies
\[
 \widehat Z_s(Qr)=Q^{s-1}\widehat Z_s(r),\qquad
 \widehat Z_t(-Pr)=P^{t-1}\widehat Z_t(-r).
\]
The phase is $\exp(2\pi i r(Qa-Pb))=\exp(-2\pi i rc)$. This is exactly
the Fourier series for the right side of~\eqref{dil:factor}. Analytic
continuation proves the claims.
\end{proof}

\subsection{A finite polynomial for every scale and derivative}
For a periodic distribution $T$, let $\mathsf P_{p,a}T$ be its pullback
under $x\mapsto px+a$. For functions it is $T(\{px+a\})$. Put
\[
 \Delta_{p,a}=\mathsf P_{p,a}\delta_0
   =\frac1p\sum_{x_0:\,px_0+a=0\ (\bmod 1)}\delta_{x_0}.
\]
Let $G_n$ be the canonical periodic finite part of $\gamma_n$, characterized
by
\[
 \langle G_n,\phi\rangle
 =\operatorname{CT}_{\epsilon\downarrow0}
      \int_\epsilon^1\gamma_n(x)\phi(x)\,dx.
\]
At each singular point $x_0$ of an affine factor use the right coordinate
$t=x-x_0$. Define
\[
 H_{p,a;n,k}=\operatorname{FP}_x
       \bigl[p^k\gamma_n^{(k)}(\{px+a\})\bigr],\qquad n,k\geq0,
\]
where $\operatorname{FP}_x$ retains the constant term of the cutoff
integral after deleting all negative powers of the local cutoffs and all
nonconstant powers of their logarithms. Cutoffs at different points tend
to zero independently. This specifies the coordinate $t$, rather than
the rescaled coordinate $pt$.

Set
\begin{align}
 E_k(z)&=\prod_{j=1}^k(1-z/j),\qquad E_0(z)=1,\notag\\
 b_{n,k}(L)&=n![z^{n+1}]e^{Lz}E_k(z)\label{dil:bdef}\\
 &=n!\sum_{j=0}^{\min(k,n+1)}(-1)^j e_j(k)
                  \frac{L^{n+1-j}}{(n+1-j)!},\notag
\end{align}
where $e_j(k)$ is the elementary symmetric polynomial of degree $j$ in
$1,1/2,\ldots,1/k$ and $e_0(k)=1$.

\begin{theorem}[Dilation and differentiation contact polynomial]
\label{dil:contact}
For all $p\geq1$, $n,k\geq0$, and real shifts $a$,
\begin{equation}\label{dil:contactformula}
 D^k\mathsf P_{p,a}G_n
 =H_{p,a;n,k}+b_{n,k}(\log p)\Delta_{p,a}^{(k)}.
\end{equation}
Consequently,
\begin{equation}\label{dil:commutator}
 D^kH_{p,a;n,0}-H_{p,a;n,k}
 =\left[b_{n,k}(\log p)-\frac{\log^{n+1}p}{n+1}\right]
       \Delta_{p,a}^{(k)}.
\end{equation}
In particular,
\begin{equation}\label{dil:scaleonly}
 \mathsf P_{p,a}G_n
 =H_{p,a;n,0}+\frac{\log^{n+1}p}{n+1}\Delta_{p,a}.
\end{equation}
\end{theorem}

\begin{proof}
Use the meromorphic periodic family
$Z(1-z)=(\delta_0-1)/z+
\sum_{n\geq0}G_nz^n/n!$. Distributional differentiation and pullback
commute with analytic continuation; for instance, this is first checked
where all endpoint derivatives vanish and then continued. Near a singular
point the $k$th argument derivative of the singular part is
\[
 (-1)^k p^{z-1}(1-z)_k\,t^{z-k-1}.
\]
Subtract the Taylor polynomial of a test function through degree $k$.
Only its degree-$k$ term produces a pole at $z=0$. In distribution
notation its contribution is
\[
 \frac{e^{z\log p}E_k(z)}{z}\,\Delta_{p,a}^{(k)}.
\]
The regular coefficient of this expression exceeds the fixed-coordinate
constant term of the coefficientwise cutoff integral by
$[z^{n+1}]e^{z\log p}E_k(z)\Delta_{p,a}^{(k)}$. The remaining Taylor
terms and the Taylor remainder are holomorphic and their coefficients
are precisely the corresponding finite parts. Multiplication by $n!$
proves~\eqref{dil:contactformula}. Its $k=0$ case gives
\eqref{dil:scaleonly}; differentiating that case and subtracting the
general identity proves~\eqref{dil:commutator}.
\end{proof}

The coefficients are explicit harmonic polynomials. For example,
\[
 b_{0,k}(L)=L-H_k,\qquad
 b_{1,k}(L)=\frac{L^2}{2}-H_kL+
                \frac{H_k^2-H_k^{(2)}}2.
\]
At $p=1$ this recovers the existing contact coefficient
$(-1)^{n+1}n!e_{n+1}(k)$. At a nonunit scale the new logarithmic terms
are essential. For example,
\[
 DH_{p,a;1,0}-H_{p,a;1,1}=-(\log p)\Delta_{p,a}',
\]
although the same commutator vanishes at $p=1$.

For completeness, define the transfer operator by
\[
 (\mathsf L_r f)(x)=\frac1r\sum_{j=0}^{r-1}f((x+j)/r),
 \qquad \widehat{\mathsf L_rT}(k)=\widehat T(rk).
\]
Its exact distribution relation is
\begin{equation}\label{dil:transfer}
 \mathsf L_rG_n=\sum_{j=0}^n\binom nj(-\log r)^{n-j}G_j
       +\frac{(-\log r)^{n+1}}{n+1}(\delta_0-1).
\end{equation}
Indeed, $\mathsf L_rZ(1-z)=r^{-z}Z(1-z)$ follows from the Fourier
coefficients, and the Laurent expansion gives~\eqref{dil:transfer}.
The operators themselves satisfy
\[
 D^k\mathsf P_{r,0}=r^k\mathsf P_{r,0}D^k,\qquad
 D^k\mathsf L_r=r^{-k}\mathsf L_rD^k,\qquad
 \mathsf L_r\mathsf P_{r,0}=\mathrm{Id}.
\]
Thus all discrepancies arise from converting the analytic distributions
to a fixed-coordinate finite part; the transfer and pullback operators
have the usual exact chain rules.

% Self-contained base closure, using lambda = 1 - s throughout.
% Intended immediately before the all-index dilation subsection.
\subsection{A self-contained Gamma quotient for the base correlation}
\label{base:subsection}

The reduction below is already present, in equivalent conventions, in the
incoming bilinear reports.  We include its short proof so that the
integer-dilation identities require no external all-order closure theorem.
For $0<c<1$, set $b=1-c$ and write
\[
 \mathcal G(z;c)=\zeta(1-z,c)+\frac1z
       =\sum_{m\ge0}\gamma_m(c)\frac{z^m}{m!}.
\]
The singularity at $z=0$ is removable.  Define the holomorphic Gamma
quotient and its removable divided difference by
\begin{align}
 U(z,w)&=\frac{\Gamma(1+z)\Gamma(1-z-w)}{\Gamma(1-w)},
 \label{base:U}\\
 V(z,w)&=\frac{U(z,w)-1}{z(z+w)}.
 \label{base:V}
\end{align}
These formulas are used near $(z,w)=(0,0)$.  Indeed, $U(0,w)=1$ and
$U(z,-z)=1$, so the numerator is divisible by both distinct linear
factors and $V$ is holomorphic.  Its exact coefficient generator is
\begin{equation}\label{base:Ugenerator}
 U(z,w)=\exp\left(
  \sum_{k=2}^{\infty}\frac{\zeta(k)}k
     \bigl[(-1)^kz^k+(z+w)^k-w^k\bigr]\right).
\end{equation}
In particular $V(0,0)=\zeta(2)$, and every Taylor coefficient is a finite
polynomial with rational coefficients in positive integer zeta values.

Let $J_{mn}(c)$ be the separated-endpoint, unscaled-coordinate finite
part of $\gamma_m(x)\gamma_n(\{x+c\})$, and set
\[
 \mathcal J(z,w;c)=\sum_{m,n\ge0}
                    J_{mn}(c)\frac{z^mw^n}{m!n!}.
\]

\begin{theorem}[Holomorphic base generating identity]
\label{base:theorem}
With $r=z+w$, one has
\begin{align}
 \mathcal J(z,w;c)={}&
 \frac{\mathcal G(r;c)-\mathcal G(w;c)}z
 +\frac{\mathcal G(r;b)-\mathcal G(z;b)}w\notag\\
 &-V(z,w)-V(w,z)\notag\\
 &+r\bigl[V(z,w)\mathcal G(r;c)
                  +V(w,z)\mathcal G(r;b)\bigr].
 \label{base:Jgenerator}
\end{align}
Every apparent singularity on the right is removable.  For $M=m+n$,
\begin{equation}\label{base:topandmissing}
 J_{mn}(c)=\frac{\gamma_{M+1}(c)}{m+1}
          +\frac{\gamma_{M+1}(1-c)}{n+1}
          +\sum_{j=0}^{M-1}
               \bigl(a_j\gamma_j(c)+b_j\gamma_j(1-c)\bigr)+d,
\end{equation}
where the sum is empty for $M=0$ and the coefficients lie in
$\mathbb Q[\zeta(2),\ldots,\zeta(M+2)]$.  Thus the index $M$ is absent.
The complete coefficients are obtained by multiplying
$[z^mw^n]$ of~\eqref{base:Jgenerator} by $m!n!$.
\end{theorem}

\begin{proof}
For $\Re z,\Re w>0$, the shifted ordinary Hurwitz integral has the
meromorphic beta representation
\begin{equation}\label{base:Cbeta}
 C(1-z,1-w;c)=B(z,1-r)\zeta(1-r,c)
                  +B(w,1-r)\zeta(1-r,b).
\end{equation}
Here the sum is continued before any exceptional parameter is evaluated.
For completeness, the Fourier calculation first holds for
$\Re z,\Re w>1$.  The positive and negative Fourier frequencies give
\[
 \frac{\Gamma(z)\Gamma(w)}{(2\pi)^r}
 \left[e^{i\pi(z-w)/2}\operatorname{Li}_r(e^{2\pi ic})
       +e^{-i\pi(z-w)/2}\operatorname{Li}_r(e^{-2\pi ic})\right].
\]
Inserting the Hurwitz Fourier formula into the right side of
\eqref{base:Cbeta}, and using the Gamma reflection identity, gives this
same expression.  The relevant phase identity is
\[
 \frac{\sin(\pi w)e^{-i\pi r/2}
             +\sin(\pi z)e^{i\pi r/2}}{\sin(\pi r)}
       =e^{i\pi(z-w)/2}.
\]
Local endpoint domination and analytic continuation establish
\eqref{base:Cbeta} meromorphically.

The unscaled endpoint finite parts give the Laurent decomposition
\begin{equation}\label{base:Laurent}
 C(1-z,1-w;c)=-\frac1{zw}
    +\frac{\mathcal G(w;c)}z
    +\frac{\mathcal G(z;b)}w+\mathcal J(z,w;c).
\end{equation}
One way to see the residue and the regular part explicitly is to use
the periodic expansion
$Z(1-z)=(\delta_0-1)/z+\sum G_mz^m/m!$.
The two shifted singular supports are disjoint; their polar product
has integral $-1$, their mixed products evaluate the smooth opposing
factor, and their regular product is precisely the prescribed local
finite part.  Equivalently, split the ordinary integral at $1-c$ and
subtract the two leading endpoint terms before expanding.  This proves
\eqref{base:Laurent} without multiplying distributions with common
singular support.

Now $B(z,1-r)=U(z,w)/z$ and
$U(z,w)=1+zrV(z,w)$.  Substitute
$\zeta(1-r,c)=-1/r+\mathcal G(r;c)$ and its reflected analogue into
\eqref{base:Cbeta}.  The two leading polar terms combine as
$-(1/z+1/w)/r=-1/(zw)$.  Subtracting the remaining polar terms in
\eqref{base:Laurent} gives exactly~\eqref{base:Jgenerator}.

The two divided differences contribute respectively
$\gamma_{M+1}(c)/(m+1)$ and $\gamma_{M+1}(b)/(n+1)$.
In the last line of~\eqref{base:Jgenerator}, the factor $r$ forces every
Stieltjes index to be at most $M-1$.  The middle line contributes only
zeta-polynomial constants.  Formula~\eqref{base:Ugenerator} then proves
the coefficient ring and the missing-index assertion.
\end{proof}

The first row is
\[
 J_{00}(c)=\gamma_1(c)+\gamma_1(1-c)-2\zeta(2).
\]
This checks both the Laurent signs and the finite-part constant.

\subsection{All Stieltjes indices in an explicit finite formula}
Define the base finite part
\[
 J_{mn}(c)=\operatorname{FP}\int_0^1
               \gamma_m(x)\gamma_n(\{x+c\})\,dx
\]
using the unscaled right coordinate at its two endpoints. Theorem~\ref{base:theorem} gives its all-order linear closure. For a
transparent finite formula put
\begin{align*}
 A_n(L;x)&=\sum_{j=0}^n\binom nj(-L)^{n-j}\gamma_j(x),\\
 u_n(L)&=\frac{(-L)^{n+1}}{n+1},\qquad
 v_n(L)=\frac{L^{n+1}}{n+1},\qquad
 L_P=\log P,\quad L_Q=\log Q.
\end{align*}

\begin{theorem}[Unequal-dilation finite-part closure]
\label{dil:allindices}
Under the separated-endpoint condition $0<c<1$, let
\[
 K_{mn}=\int_{\mathbb T}(\mathsf P_{p,a}G_m)
                        (\mathsf P_{q,b}G_n).
\]
This product is well-defined because its factors have disjoint singular
supports. Then
\begin{align}
 K_{mn}={}&\sum_{i=0}^m\sum_{j=0}^n
   \binom mi\binom nj(-L_Q)^{m-i}(-L_P)^{n-j}J_{ij}(c)\notag\\
 &+u_m(L_Q)A_n(L_P;c)+u_n(L_P)A_m(L_Q;1-c)\notag\\
 &-u_m(L_Q)u_n(L_P).\label{dil:pullformula}
\end{align}
The finite part using the coordinate $x-x_0$ at every singular point is
\begin{align}
 \operatorname{FP}_x\int_0^1
     \gamma_m(\{px+a\})\gamma_n(\{qx+b\})\,dx
 ={}&K_{mn}-v_m(\log p)\,[A_n(L_P;c)-u_n(L_P)]\notag\\
 &-v_n(\log q)\,[A_m(L_Q;1-c)-u_m(L_Q)].
 \label{dil:intrinsicformula}
\end{align}
Thus every index pair still reduces to a finite linear combination of
$\gamma_j(c)$ and $\gamma_j(1-c)$ for $0\leq j\leq m+n+1$, with
coefficients in the ring generated over $\mathbb Q$ by positive integer
zeta values and $\log p,\log q,\log P,\log Q$.
\end{theorem}

\begin{proof}
Write $\mathcal G(z;x)=\sum_{n\geq0}\gamma_n(x)z^n/n!$ and
$\mathcal J(z,w;c)=\sum_{m,n\geq0}J_{mn}(c)z^mw^n/(m!n!)$.
The separated-endpoint Laurent expansion is
\[
 C(1-z,1-w;c)=-\frac1{zw}
       +\frac{\mathcal G(w;c)}z
       +\frac{\mathcal G(z;1-c)}w+\mathcal J(z,w;c).
\]
Multiply this by $Q^{-z}P^{-w}$ as required by
Theorem~\ref{dil:master}, and take the regular coefficient $m!n![z^mw^n]$.
The four displayed terms give the four terms in
\eqref{dil:pullformula}. This coefficient is the product of the regular
coefficients of the two pulled-back distribution families.

To change the endpoint coordinates, apply~\eqref{dil:scaleonly} to the
two factors. Their delta combs have disjoint supports, so the product of
the two combs is zero and the two remaining correction terms evaluate
the smooth opposing factor. At the first comb,
\[
 \frac1p\sum_{px_0+a\in\mathbb Z}\gamma_n(\{qx_0+b\})
 =\frac1P\sum_{r=0}^{P-1}\gamma_n((r+c)/P)
 =A_n(L_P;c)-u_n(L_P).
\]
The first equality uses coprimality of $P,Q$; the second follows by
expanding the Hurwitz multiplication formula at $s=1$. The second comb
gives $A_m(L_Q;1-c)-u_m(L_Q)$. This proves
\eqref{dil:intrinsicformula}. The last assertion follows from the base
linear closure and the explicitly triangular index sums.
\end{proof}

At the first Stieltjes index, the existing identity
$J_{00}(c)=\gamma_1(c)+\gamma_1(1-c)-2\zeta(2)$ gives the particularly
short pulled-back answer
\[
 K_{00}=\gamma_1(c)+\gamma_1(1-c)-2\zeta(2)
       -L_Q\gamma_0(c)-L_P\gamma_0(1-c)-L_PL_Q.
\]
For $p=1,q=2$, the fixed-$x$ version is
\begin{align*}
 \operatorname{FP}_x\int_0^1\psi(x)\psi(\{2x+c\})\,dx
 ={}&\gamma_1(c)+\gamma_1(1-c)-2\zeta(2)\\
 &+(\log2)[\psi(c)+\psi(1-c)]-\log^22.
\end{align*}
Replacing a pair $(p,q)$ by $(dp,dq)$ leaves the ordinary spectral
integral unchanged, but in general changes its fixed-coordinate finite
part. Formula~\eqref{dil:intrinsicformula} gives the complete discrepancy.

\subsection{An ordinary log-Gamma identity with unequal frequencies}
Let $\kappa=\gamma+\log(2\pi)$, $z_c=e^{2\pi ic}$, and let
$\partial_\nu$ mean differentiation of the polylogarithm with respect
to its order. The following identity needs no finite part.

\begin{corollary}[Integer-affine log-Gamma correlation]
\label{dil:loggamma}
For the same $p,q,a,b$ and $0<c<1$,
\begin{align}
 &\int_0^1\log\Gamma(\{px+a\})\log\Gamma(\{qx+b\})\,dx
   =\frac{\log^2(2\pi)}4\notag\\
 &\quad+\frac1{2\pi^2PQ}\operatorname{Re}
  \left[\left(\partial_\nu^2-(2\kappa+L_P+L_Q)\partial_\nu
       +(\kappa+L_P)(\kappa+L_Q)+\frac{\pi^2}4\right)
           \operatorname{Li}_\nu(z_c)\right]_{\nu=2}\notag\\
 &\quad+\frac{L_P-L_Q}{4\pi PQ}
           \operatorname{Im}\operatorname{Li}_2(z_c).
 \label{dil:loggammaformula}
\end{align}
The formula extends continuously to $c=0$ for this ordinary integral.
\end{corollary}

\begin{proof}
Lerch's identity gives the mean-zero function
$L(x)=\zeta'(0,x)=\log\Gamma(x)-\tfrac12\log(2\pi)$. Its nonzero
Fourier coefficients are
\[
 \widehat L(r)=\frac{\kappa+\log|r|+i\pi\operatorname{sgn}(r)/2}
                       {2\pi i r}.
\]
Either differentiate~\eqref{dil:factor} at $(s,t)=(0,0)$ or apply
Parseval to these square-summable coefficients. Positive and negative
frequencies yield
\begin{align*}
 &\int_0^1L(\{px+a\})L(\{qx+b\})\,dx\\
 &\quad=\frac1{2\pi^2PQ}\sum_{r=1}^\infty\frac1{r^2}
 \bigg[\left((\kappa+L_Q+\log r)(\kappa+L_P+\log r)
          +\frac{\pi^2}4\right)\cos(2\pi rc)\\
 &\hspace{30mm}+\frac\pi2(L_P-L_Q)\sin(2\pi rc)\bigg].
\end{align*}
The series is absolutely convergent, including two order derivatives,
and converts termwise to~\eqref{dil:loggammaformula}. Restoring the mean
gives the constant term. Absolute convergence also proves continuity at
$c=0$.
\end{proof}

\subsection{Audit finding and scope}
The incoming archive \nolinkurl{ProveIt_Stieltjes_Convolution_2026-10-10.zip},
\texttt{article.tex}, lines 1118--1121, suggests that changing one affine
frequency produces multiple independent Fourier indices. That statement
needs qualification: for two nonzero integer frequencies, the equation
$pn+qm=0$ always has the one-index parametrization used above. A suitable
replacement is: ``Two periodic factors with positive integer affine
frequencies retain a one-dimensional frequency constraint. Their complete
shifted reduction, with coordinate-dependent finite parts, is given by
the integer-dilation theorem. Products of three or more independent
factors generally leave a higher-dimensional frequency lattice.'' This
corrects an explanatory sentence in a future-work section; it does not
invalidate the convolution theorems of that report.

The new formulas answer the separated-endpoint unequal-dilation problem
posed in \nolinkurl{ProveIt_Stieltjes_Correlation_Closure},
\texttt{article.tex}, lines 1247--1254, and the scale/derivative commutator
problem posed in \nolinkurl{ProveIt_Shifted_Hurwitz_Jets},
\texttt{sections.tex}, lines 782--783. They do not settle collision
renormalization for two singularities at the same point or prove
arithmetic independence of the resulting constant coordinates.
